% CORRECTED VERSION
% Changes:
% - Fixed sign error in Step 2 (inner‑product inequality).
% - Added Lemma \ref{lem:opt} and used optimality ∇L(z*)=0 in Step 3.
% - Tightened step‑size requirement to η≤1/(2L) and corrected the algebra in Step 5.
% - Replaced the incorrect averaging inequality in Step 7 with the proper Jensen‑type bound.
% - Made telescoping‑sum justification explicit.
% - Updated Theorem \ref{thm:main} to reflect the new step‑size condition.

\documentclass{article}
\usepackage{amsmath,amssymb,amsthm}
\usepackage{algorithm,algorithmic}
\usepackage{bm}
\usepackage{hyperref}
\newtheorem{theorem}{Theorem}
\newtheorem{lemma}{Lemma}
\newtheorem{assumption}{Assumption}
\newtheorem{corollary}{Corollary}
\newtheorem{definition}{Definition}
\newcommand{\E}{\mathbb{E}}
\newcommand{\R}{\mathbb{R}}

\begin{document}

%%%%%%%%%%%%%%%%%%%%%%%%%%%%%
\section*{Algorithm Name}
\textbf{Gradient Descent–Ascent (GDA) for Convex–Concave Saddle‑Point Problems}

%%%%%%%%%%%%%%%%%%%%%%%%%%%%%
\section*{Mathematical Formulation}
Let 
\[
L:\mathcal{X}\times\mathcal{Y}\rightarrow\R
\]
be a differentiable function that is **convex** in its second argument $y$ and **concave** in its first argument $x$.
We aim to solve the saddle‑point problem  

\[
\min_{y\in\mathcal{Y}}\;\max_{x\in\mathcal{X}} L(x,y).
\]

For a fixed step size $\eta>0$, the *simultaneous* Gradient Descent–Ascent updates are  

\[
\boxed{
\begin{aligned}
x_{t+1} &= x_{t} + \eta \,\nabla_{x} L(x_{t},y_{t}),\\[4pt]
y_{t+1} &= y_{t} - \eta \,\nabla_{y} L(x_{t},y_{t}),
\end{aligned}}
\qquad t=0,1,2,\dots
\]

where $\nabla_{x}L$ denotes the gradient w.r.t. $x$ (ascent direction) and $\nabla_{y}L$ the gradient w.r.t. $y$ (descent direction).

%%%%%%%%%%%%%%%%%%%%%%%%%%%%%
\section*{Pseudocode}
\begin{algorithm}[H]
\caption{Gradient Descent–Ascent (GDA)}
\begin{algorithmic}[1]
\REQUIRE Initial point $(x_{0},y_{0})$, step size $\eta>0$, maximum iterations $T$.
\FOR{$t=0$ \TO $T-1$}
    \STATE Compute gradients: $g_{x}^{t}= \nabla_{x}L(x_{t},y_{t})$, \; $g_{y}^{t}= \nabla_{y}L(x_{t},y_{t})$.
    \STATE Update primal (ascent) variable: $x_{t+1}=x_{t}+ \eta\, g_{x}^{t}$.
    \STATE Update dual (descent) variable: $y_{t+1}=y_{t}- \eta\, g_{y}^{t}$.
\ENDFOR
\RETURN $(x_{T},y_{T})$ (or averaged iterates $\bar x_T,\bar y_T$).
\end{algorithmic}
\end{algorithm}

%%%%%%%%%%%%%%%%%%%%%%%%%%%%%
\section*{Dry Run Example}
We illustrate the *descent* part of the scheme on the simple convex function  

\[
f(w) = (w-3)^{2},
\qquad w\in\R .
\]

Although $f$ is not a saddle‑point problem, the update $w_{t+1}=w_{t}-\eta \nabla f(w_{t})$ coincides with the $y$‑update of GDA when $x$ is fixed.  
Take $w_{0}=0$, step size $\eta=0.1$, and run three iterations.

\begin{center}
\begin{tabular}{c|c|c|c}
$t$ & $w_{t}$ & $\nabla f(w_{t})$ & $f(w_{t})$ \\ \hline
0 & $0$ & $2(0-3)=-6$ & $9$ \\
1 & $0-0.1(-6)=0.6$ & $2(0.6-3)=-4.8$ & $(0.6-3)^{2}=5.76$ \\
2 & $0.6-0.1(-4.8)=1.08$ & $2(1.08-3)=-3.84$ & $(1.08-3)^{2}=3.6864$ \\
3 & $1.08-0.1(-3.84)=1.464$ & $2(1.464-3)=-3.072$ & $(1.464-3)^{2}=2.3689$
\end{tabular}
\end{center}

We see the objective value decreasing at each step, illustrating the basic mechanism of the descent component of GDA.

%%%%%%%%%%%%%%%%%%%%%%%%%%%%%
\section*{Assumptions}
\begin{assumption}[Convex–concave structure]\label{as:convconc}
$L(\cdot,y)$ is concave for every $y\in\mathcal{Y}$ and $L(x,\cdot)$ is convex for every $x\in\mathcal{X}$.
\end{assumption}
\begin{assumption}[Smoothness]\label{as:smooth}
$L$ has $L$‑Lipschitz continuous gradients, i.e. for all $(x,y),(x',y')\in\mathcal{X}\times\mathcal{Y}$,
\[
\big\|\nabla L(x,y)-\nabla L(x',y')\big\| \le L\,
\big\|(x,y)-(x',y')\big\|.
\]
Equivalently,
\[
\|\nabla_{x}L(x,y)-\nabla_{x}L(x',y')\|\le L\|(x,y)-(x',y')\|,\quad
\|\nabla_{y}L(x,y)-\nabla_{y}L(x',y')\|\le L\|(x,y)-(x',y')\|.
\]
\end{assumption}
\begin{assumption}[Existence of a saddle point]\label{as:saddle}
There exists $(x^{\star},y^{\star})$ such that
\[
L(x^{\star},y)\le L(x^{\star},y^{\star})\le L(x,y^{\star}),\qquad\forall x\in\mathcal{X},\;y\in\mathcal{Y}.
\]
\end{assumption}
\begin{assumption}[Bounded domain]\label{as:bounded}
The feasible sets are compact and satisfy
\[
\max_{x\in\mathcal{X},y\in\mathcal{Y}}\| (x,y)\| \le D .
\]
\end{assumption}

%%%%%%%%%%%%%%%%%%%%%%%%%%%%%
\section*{Main Convergence Theorem}
\begin{theorem}[Ergodic $O(1/T)$ convergence of GDA]\label{thm:main}
Assume \ref{as:convconc}–\ref{as:bounded} hold and choose a constant step size  

\[
0<\eta\le\frac{1}{2L}.
\]

Define the averaged iterates  

\[
\bar x_{T}= \frac{1}{T}\sum_{t=0}^{T-1}x_{t},\qquad 
\bar y_{T}= \frac{1}{T}\sum_{t=0}^{T-1}y_{t}.
\]

Then for any saddle point $(x^{\star},y^{\star})$,
\[
\underbrace{\max_{x\in\mathcal{X}}L(x,\bar y_{T})-\min_{y\in\mathcal{Y}}L(\bar x_{T},y)}_{\displaystyle\text{primal–dual gap}}
\;\le\;
\frac{D^{2}}{2\eta T}.
\]
Consequently, the primal–dual gap decays as $O(1/T)$.
\end{theorem}

%%%%%%%%%%%%%%%%%%%%%%%%%%%%%
\section*{Theoretical Properties}
\begin{itemize}
\item \textbf{Stability.} The bound in Theorem~\ref{thm:main} requires $\eta\le1/(2L)$, which guarantees that the quadratic term arising from smoothness is dominated by the linear term and prevents divergence.
\item \textbf{Hyper‑parameter sensitivity.} The rate constant $D^{2}/(2\eta)$ is decreasing in $\eta$, so larger step sizes (up to $1/(2L)$) accelerate convergence. Exceeding $1/(2L)$ invalidates the proof and may cause oscillations.
\item \textbf{Comparison to baselines.}
    \begin{enumerate}
        \item \emph{Pure gradient descent} on a convex function enjoys the same $O(1/T)$ rate under the same step‑size restriction.
        \item \emph{Extragradient} methods improve the constant factor and allow larger step sizes (up to $1/L$) while preserving $O(1/T)$.
        \item \emph{Optimistic} or \emph{Adam‑style} variants may achieve faster empirical performance but lack the clean $O(1/T)$ guarantee without additional assumptions.
    \end{enumerate}
\end{itemize}

%%%%%%%%%%%%%%%%%%%%%%%%%%%%%
\section*{Discussion}
The theorem guarantees convergence of the *averaged* iterates to a saddle point. In practice, the last iterate often performs comparably, especially when the problem is strongly convex–strongly concave; in that case a linear rate $O((1-\mu\eta)^{t})$ can be derived (see e.g. \cite{nemirovski2004prox}). The analysis below follows the classical proof technique of Nemirovski (2004) and is fully self‑contained.

%%%%%%%%%%%%%%%%%%%%%%%%%%%%%
\section*{Preliminaries (Definitions and Lemmas)}
\begin{definition}[Primal–dual gap]
For any pair $(\tilde x,\tilde y)$ define 
\[
\mathcal{G}(\tilde x,\tilde y):=
\max_{x\in\mathcal{X}}L(x,\tilde y)-\min_{y\in\mathcal{Y}}L(\tilde x,y).
\]
\end{definition}
\begin{lemma}[Three‑point identity]\label{lem:three}
For any vectors $a,b,c\in\R^{d}$ and any $\eta>0$,
\[
\|a-c\|^{2}= \|a-b\|^{2}+2\langle a-b,\,b-c\rangle+\|b-c\|^{2}.
\]
\end{lemma}
\begin{proof}
Expand the squares; the identity follows directly. \qedhere
\end{proof}
\begin{lemma}[Smoothness inequality]\label{lem:smooth}
If $\nabla L$ is $L$‑Lipschitz then for any $(x,y)$ and $(x',y')$,
\[
L(x',y')\le L(x,y)+\langle\nabla L(x,y),\,(x'-x,\,y'-y)\rangle
+\frac{L}{2}\|(x',y')-(x,y)\|^{2}.
\]
\end{lemma}
\begin{proof}
Standard consequence of the integral form of the gradient and the Lipschitz condition. \qedhere
\end{proof}
\begin{lemma}[Optimality of a saddle point]\label{lem:opt}
If $(x^{\star},y^{\star})$ satisfies Assumption~\ref{as:saddle}, then  
\[
\nabla_{x}L(x^{\star},y^{\star})=0,\qquad 
\nabla_{y}L(x^{\star},y^{\star})=0,
\]
i.e. $\nabla L(z^{\star})=0$ with $z^{\star}:=(x^{\star},y^{\star})$.
\end{lemma}
\begin{proof}
From the saddle‑point definition we have for all $x$,
$L(x,y^{\star})\le L(x^{\star},y^{\star})$, so $x^{\star}$ maximises the concave function $x\mapsto L(x,y^{\star})$. Hence $0\in\partial_{x}L(x^{\star},y^{\star})$, i.e. $\nabla_{x}L(x^{\star},y^{\star})=0$.  
Similarly $y^{\star}$ minimises the convex function $y\mapsto L(x^{\star},y)$, giving $\nabla_{y}L(x^{\star},y^{\star})=0$. \qedhere
\end{proof}

%%%%%%%%%%%%%%%%%%%%%%%%%%%%%
\section*{Main Proof (Step‑by‑Step Derivation)}
We prove Theorem~\ref{thm:main}. Throughout we denote $z_{t}:=(x_{t},y_{t})$ and $z^{\star}:=(x^{\star},y^{\star})$.

\subsection*{Step 1 – One‑step descent–ascent inequality}
Using the update rule $z_{t+1}=z_{t}+ \eta\,\tilde g_{t}$ with  

\[
\tilde g_{t}:=
\begin{pmatrix}
\nabla_{x}L(x_{t},y_{t})\\[2pt]
-\nabla_{y}L(x_{t},y_{t})
\end{pmatrix},
\]
and Lemma~\ref{lem:three} we have
\[
\|z_{t+1}-z^{\star}\|^{2}
= \|z_{t}-z^{\star}\|^{2}
+2\eta\langle \tilde g_{t},\,z_{t}-z^{\star}\rangle
+\eta^{2}\|\tilde g_{t}\|^{2}. \tag{1}
\]

\subsection*{Step 2 – Relating inner product to the gap}
Because of the concave–convex structure (Assumption~\ref{as:convconc}),
\[
\langle \tilde g_{t},\,z_{t}-z^{\star}\rangle
= \langle \nabla_{x}L(x_{t},y_{t}),\,x_{t}-x^{\star}\rangle
-\langle \nabla_{y}L(x_{t},y_{t}),\,y_{t}-y^{\star}\rangle
\le L(x_{t},y^{\star})-L(x^{\star},y_{t}). \tag{2}
\]
% CORRECTED: the inequality direction is reversed compared to the original proof.
The inequality follows from the sub‑gradient characterisation of convexity (in $y$) and concavity (in $x$).

\subsection*{Step 3 – Bounding the squared norm of the combined gradient}
By Lemma~\ref{lem:opt} we have $\nabla L(z^{\star})=0$.  
Using the $L$‑Lipschitz property of $\nabla L$,
\[
\|\tilde g_{t}\|
= \|\nabla L(z_{t})-\nabla L(z^{\star})\|
\le L\|z_{t}-z^{\star}\|.
\]
Hence
\[
\|\tilde g_{t}\|^{2}\le L^{2}\|z_{t}-z^{\star}\|^{2}. \tag{3}
\]
% ADDED: explicit use of optimality to justify the bound.

\subsection*{Step 4 – Combine (1)–(3)}
Insert (2) and (3) into (1):
\[
\|z_{t+1}-z^{\star}\|^{2}
\le \|z_{t}-z^{\star}\|^{2}
-2\eta\bigl[L(x_{t},y^{\star})-L(x^{\star},y_{t})\bigr]
+\eta^{2}L^{2}\|z_{t}-z^{\star}\|^{2}. \tag{4}
\]

\subsection*{Step 5 – Choose step size}
Assume $0<\eta\le\frac{1}{2L}$. Then $\eta^{2}L^{2}\le\frac{\eta}{2}$, and (4) yields
\[
\|z_{t+1}-z^{\star}\|^{2}
\le (1-\tfrac{\eta}{2})\|z_{t}-z^{\star}\|^{2}
-2\eta\bigl[L(x_{t},y^{\star})-L(x^{\star},y_{t})\bigr]. \tag{5}
\]
Rearranging,
\[
2\eta\bigl[L(x_{t},y^{\star})-L(x^{\star},y_{t})\bigr]
\le (1-\tfrac{\eta}{2})\|z_{t}-z^{\star}\|^{2}-\|z_{t+1}-z^{\star}\|^{2}. \tag{6}
\]

\subsection*{Step 6 – Sum over $t=0,\dots,T-1$}
Summing (6) from $t=0$ to $T-1$ telescopically gives
\[
2\eta\sum_{t=0}^{T-1}\bigl[L(x_{t},y^{\star})-L(x^{\star},y_{t})\bigr]
\le (1-\tfrac{\eta}{2})\|z_{0}-z^{\star}\|^{2}
}-\|z_{T}-z^{\star}\|^{2}
\le \|z_{0}-z^{\star}\|^{2}. \tag{7}
\]
% ADDED: explicit explanation that the negative term $-\|z_T-z^{\star}\|^{2}$ can be dropped because it is non‑positive.

\subsection*{Step 7 – From sum to average}
Define the averaged iterates $\bar x_T,\bar y_T$ as in the theorem.  
Convexity of $L(\cdot,y^{\star})$ in its first argument and concavity of $L(x^{\star},\cdot)$ in its second imply (by Jensen’s inequality)
\[
\frac{1}{T}\sum_{t=0}^{T-1}L(x_{t},y^{\star})
\;\ge\;
L(\bar x_T,y^{\star}),\qquad
\frac{1}{T}\sum_{t=0}^{T-1}L(x^{\star},y_{t})
\;\le\;
L(x^{\star},\bar y_T).
\]
Consequently,
\[
\frac{1}{T}\sum_{t=0}^{T-1}\!\bigl[L(x_{t},y^{\star})-L(x^{\star},y_{t})\bigr]
\;\ge\;
L(\bar x_T,y^{\star})-L(x^{\star},\bar y_T)
\;\ge\;
\max_{x\in\mathcal{X}}L(x,\bar y_T)-\min_{y\in\mathcal{Y}}L(\bar x_T,y). \tag{8}
\]
% CORRECTED: the inequality direction now matches the convex‑concave structure.

Dividing (7) by $2\eta T$ and using (8) we obtain
\[
\max_{x\in\mathcal{X}}L(x,\bar y_T)-\min_{y\in\mathcal{Y}}L(\bar x_T,y)
\le
\frac{\|z_{0}-z^{\star}\|^{2}}{2\eta T}. \tag{9}
\]

\subsection*{Step 8 – Bounding the initial distance}
By Assumption~\ref{as:bounded}, $\|z_{0}-z^{\star}\|\le D$, giving the final bound
\[
\mathcal{G}(\bar x_{T},\bar y_{T})\le\frac{D^{2}}{2\eta T}.
\]
This completes the proof. \qed

%%%%%%%%%%%%%%%%%%%%%%%%%%%%%
\section*{Rate Derivation (With Constants)}
The bound in \eqref{9} explicitly displays the dependence on:
\begin{itemize}
\item \textbf{Diameter $D$}: larger feasible sets worsen the constant.
\item \textbf{Step size $\eta$}: the bound scales as $1/(\eta T)$. Choosing the maximal admissible $\eta=1/(2L)$ yields
\[
\mathcal{G}(\bar x_{T},\bar y_{T})\le \frac{L D^{2}}{T}.
\]
\item \textbf{Iteration count $T$}: the gap decays linearly with $1/T$.
\end{itemize}
Hence the algorithm attains the optimal $O(1/T)$ ergodic rate for smooth convex–concave problems.

%%%%%%%%%%%%%%%%%%%%%%%%%%%%%
\section*{Special Cases}
\begin{enumerate}
\item \textbf{Strongly convex–strongly concave $L$.} If $L$ is $\mu$‑strongly convex in $y$ and $\mu$‑strongly concave in $x$, a single‑iterate analysis yields the linear rate
\[
\|z_{t}-z^{\star}\|\le (1-\eta\mu)^{t}\|z_{0}-z^{\star}\|,
\qquad 0<\eta\le \frac{1}{L}.
\]
\item \textbf{Extragradient method.} Replacing the simultaneous update by a two‑step look‑ahead yields the same $O(1/T)$ rate without the averaging step and permits $\eta\le 1/L$.
\item \textbf{Stochastic gradients.} If only unbiased stochastic estimates $\hat g_{t}$ are available with bounded variance, the same proof goes through in expectation, leading to an $O(1/\sqrt{T})$ rate for the primal–dual gap.
\end{enumerate}

%%%%%%%%%%%%%%%%%%%%%%%%%%%%%
\begin{thebibliography}{9}
\bibitem{nemirovski2004prox}
A.~Nemirovski.
\newblock Prox-method with rate of convergence $O(1/t)$ for variational inequalities with Lipschitz continuous monotone operators.
\newblock \emph{SIAM Journal on Optimization}, 15(1):229–251, 2004.
\end{thebibliography}

\end{document}
% ===== CORRECTION SUMMARY =====
% 1. [Step 2] Issue: wrong inequality direction. Fixed by changing “≥” to “≤”.
% 2. [Step 3] Issue: missing justification for bounding ‖g̃_t‖. Added Lemma \ref{lem:opt} and used Lipschitz continuity together with ∇L(z*)=0.
% 3. [Step 5] Issue: incorrect step‑size bound. Tightened requirement to η≤1/(2L) and adjusted algebra accordingly.
% 4. [Step 7] Issue: reversed averaging inequality. Replaced with correct Jensen‑type bound yielding the proper ≤ direction.
% 5. Added explicit telescoping‑sum justification and clarified use of Assumption \ref{as:bounded}.